\documentclass[11pt]{article}
\usepackage[margin=1in]{geometry}
\usepackage{amsmath,amssymb,amsthm}
\usepackage{booktabs}
\usepackage[colorlinks=true,linkcolor=blue,urlcolor=blue]{hyperref}

\newtheorem{claim}{Claim}
\newtheorem{remark}{Remark}

\title{A Double Counterexample to TLMC Conjecture 00000001190:\\
Prime-Order Elements of $\mathrm{GL}_2(\mathbb{F}_7)$}
\author{TLMC Submission --- Disproof of Conjecture 00000001190}
\date{September 2026}

\begin{document}
\maketitle

\begin{abstract}
We refute Conjecture 00000001190, which asserts that the number $N_p$ of
elements of prime order $p$ in $\mathrm{GL}_n(\mathbb{F}_q)$ equals
$\sum_{r \mid n,\ r \text{ prime}} \Theta_r(q)$ with
$\Theta_r(q) = q^{n^2-r}(q-1)^{-1}\prod_{i=1}^{r-1}(q^{n-i}-1)$,
``with zero error.'' At the single parameter point $(n,q) = (2,7)$ the formula
is wrong twice: it predicts $49$ elements of order $2$ where the truth is $57$,
and it predicts $0$ elements of order $3$ where the truth is $170$. Every
ground-truth number is a brute-force enumeration of all $7^4 = 2401$ matrices
over $\mathbb{Z}/7\mathbb{Z}$, independently machine-verified in Python and by a
zero-axiom Lean~4 kernel check.
\end{abstract}

\section{The conjecture}

\begin{quote}
\textit{Definition.} $N_p$ is the number of elements of prime order $p$ in
$\mathrm{GL}_n(\mathbb{F}_q)$.

\textit{Conjecture.} $N_p = \sum_{r \mid n,\ r \text{ prime}} \Theta_r(q)$,
where
\[
\Theta_r(q) = q^{n^2-r}\,(q-1)^{-1}\prod_{i=1}^{r-1}\bigl(q^{n-i}-1\bigr);
\]
this formula covers all prime-order elements with zero error.
\end{quote}

\section{The formula side at $(n,q) = (2,7)$}

Since $n = 2$, the only prime divisor of $n$ is $r = 2$. Hence the conjectured
sum has exactly one term:
\begin{align}
N_2 &\overset{!}{=} \Theta_2(7)
      = 7^{4-2}\cdot\frac{7^{2-1}-1}{7-1}
      = 49 \cdot \frac{6}{6} = \boxed{49},\\
N_3 &\overset{!}{=} 0 \quad (\text{the sum over prime } r \mid 2 \text{ contains no } r = 3,
      \text{ so it is the empty sum}),\\
N_p &\overset{!}{=} 0 \quad \text{for every prime } p \ge 3.
\end{align}

\section{Ground truth: enumeration of $\mathrm{GL}_2(\mathbb{F}_7)$}

First,
\[
\bigl|\mathrm{GL}_2(\mathbb{F}_7)\bigr| = (7^2-1)(7^2-7) = 48 \cdot 42 = 2016.
\]
This was confirmed by enumerating all $7^4 = 2401$ matrices
$\left(\begin{smallmatrix} a & b \\ c & d\end{smallmatrix}\right)$ with
$a,b,c,d \in \mathbb{Z}/7\mathbb{Z}$ and keeping exactly those with
$ad - bc \not\equiv 0 \pmod 7$; the count is $2016$. Elements of order exactly
$k$ are then those $M$ with $M \ne I$, $M^k = I$ (and automatically $M^j \ne I$
for $1 \le j < k$ since $k$ is prime). The exhaustive enumeration gives:

\begin{table}[h]
\centering
\begin{tabular}{lrr}
\toprule
quantity & brute-force enumeration & conjectured formula \\
\midrule
$|\mathrm{GL}_2(\mathbb{F}_7)|$                 & 2016 & --- \\
elements of order exactly $2$ ($N_2$)           & \textbf{57}  & \textbf{49} \\
elements of order exactly $3$ ($N_3$)           & \textbf{170} & \textbf{0} \\
\bottomrule
\end{tabular}
\caption{The conjecture fails twice at $(n,q) = (2,7)$: $57 \ne 49$ and $170 \ne 0$.}
\end{table}

As a sanity check, the full order distribution of the $2016$ elements is
$1{:}\,1$, $2{:}\,57$, $3{:}\,170$, $4{:}\,42$, $6{:}\,618$, $7{:}\,48$,
$8{:}\,84$, $12{:}\,84$, $14{:}\,48$, $16{:}\,168$, $21{:}\,96$, $24{:}\,168$,
$42{:}\,96$, $48{:}\,336$, which sums to $2016$.

\section{Why the formula undercounts}

\subsection{Order $2$: $\mathbf{57}$ against $\mathbf{49}$}

Over $\mathbb{F}_7$, $x^2 - 1 = (x-1)(x+1)$ splits completely, so every
involution is diagonalizable. The elements with eigenvalues $\{1,-1\}$
($\ne \pm I$) form a single conjugacy class: the centralizer of such an element
is the group of diagonal matrices, of order $(7-1)^2 = 36$, so the class has
size $2016/36 = 56$. Together with the scalar matrix $-I$ (which also has order
exactly $2$) this gives
\[
N_2 = 56 + 1 = 57 \ne 49 = \Theta_2(7).
\]
The term $\Theta_2(q) = q^{n^2-2}(q^1-1)/(q-1) = q^{n^2-2} = q^2$ counts the
regular-semisimple split-torus contribution (pairs of distinct eigenvalues in
$\mathbb{F}_q^\times$, modulo the multiplicative structure of the torus), but it
accounts neither for the scalar involution $-I$ nor for the fact that a
$2$-dimensional representation can realize order $2$ with eigenvalue set
$\{1,-1\}$ embedded in the conjugation structure differently than the formula
assumes.

\subsection{Order $3$: $\mathbf{170}$ against $\mathbf{0}$}

In $\mathbb{F}_7^\times$ (cyclic of order $6$) there are two primitive cube
roots of unity: $x^3 - 1 = (x-1)(x-2)(x-4) \pmod 7$, indeed
$2^3 = 8 \equiv 1$ and $4^3 = 64 \equiv 1 \pmod 7$. An element of order
exactly $3$ is diagonalizable with eigenvalues in
$\{1,2,4\}$, not both $1$. The three eigenvalue multisets
$\{2,1\}$, $\{4,1\}$, $\{2,4\}$ give three conjugacy classes, each with
centralizer of order $36$, i.e.\ $56$ elements each; plus the two scalar
matrices $2I$ and $4I$. Hence
\[
N_3 = 3 \cdot 56 + 2 = 170 \ne 0.
\]
The conjecture's divisor condition ``$r \mid n$'' is blind to this: because
$3 \mid (q-1) = 6$, the field $\mathbb{F}_7$ itself contains elements of order
$3$, and scalars $\lambda I$ of order $3$ exist in \emph{every} dimension
$n \ge 1$, as do diagonalizable ones for all $n \ge 2$. No sum of the form
$\sum_{r \mid n} \Theta_r$ can ever count them unless every prime divisor of
$q - 1$ happens to divide $n$.

\section{Machine verification}

All counts are verified by two independent, axiom-free computations:

\begin{itemize}
  \item \texttt{reproduce.py}: pure-Python enumeration of all $2401$ matrices;
        prints and asserts $2016$, $57$, $170$, $49$, $0$.
  \item \texttt{lean4/}: Lean~4 (toolchain \texttt{v4.33.1}) theorems proved by
        \texttt{decide} over the same exhaustive enumeration:
        \texttt{count\_order2 : \ldots = 57},
        \texttt{count\_order3 : \ldots = 170},
        \texttt{formula2 : $\Theta_2(7) = 49$}, and
        \texttt{refute : $57 \ne 49 \wedge 170 \ne 0$}.
        \texttt{\#print axioms} confirms every theorem
        ``\texttt{does not depend on any axioms}'' (no \texttt{sorry}, no
        \texttt{Classical.choice}, no gaps).
\end{itemize}

\section{Scope and boundaries}

We refute only the literal, universal statement of Conjecture 00000001190 at
the parameter point $(n,q) = (2,7)$, where it fails for two different primes
$p$ simultaneously. We make no claim about the behaviour of the formula at
other parameters; it may be exact or near-exact elsewhere (e.g.\ when
$\gcd(n, q-1) = 1$ the scalar-matrix obstruction disappears, though the
$57$-vs-$49$ discrepancy at $p = 2$ here suggests the semisimple counting in
$\Theta_r$ itself also needs care). Fixing the conjecture would require at
least (i) adding scalar contributions $\sum_{\lambda \in \mathbb{F}_q^\times,\
\mathrm{ord}(\lambda) = p} 1$, and (ii) handling non-regular semisimple classes
whenever $p \mid q - 1$, independently of any divisibility of $n$ by $p$.

\end{document}
